\section{代码模型生态：从 Copilot 到 BigCode 协作}

数据管线完成后，课程回到 code LLM 的历史和生态。核心变化是 Codex/Copilot 证明 next-token prediction 足以成为强 baseline，社区随后把公开数据、模型和 fine-tuning 组合成大量 code assistants。

\subsection{2021：代码可以像文本一样做 next-token prediction}

早期 code modeling 常显式使用 AST、program analysis 或 feature engineering；Codex 展示把大 Transformer 喂入大量 code，也能学会 completion。结构化方法仍有价值，但不再是进入强 code generation 的必要前置。

\lecturefigure{slide-55.jpg}{GitHub Copilot 2021：next-token code modeling 的产品起点}{Loubna Ben Allal 官方 deck，第 55 页}

\teachervoice{00:38:30--00:40:00，讲者把 Codex 的历史意义总结为“code can be treated like text”：不是代码没有结构，而是强 sequence baseline 把手工特征的门槛大幅降低。}

\subsection{2024：1.7K+ open code models 与 leaderboard landscape}

Hub 上出现大量 code-pretrained 或含 code mixture 的模型；同时 base 与 instruction-tuned models 在 HumanEval 类榜单上快速提升。数量说明 ecosystem 活跃，不等于 1,700 个独立高质量 pretraining runs，许多是 fine-tunes、quantizations 或 derivatives。

\lecturefigure{slide-56.jpg}{2024 年快照：超过 1.7K open models trained on code}{Loubna Ben Allal 官方 deck，第 56 页}

\lecturefigure{slide-57.jpg}{强 code base 与 instruction-tuned models 的榜单快照}{Loubna Ben Allal 官方 deck，第 57 页}

\begin{warningbox}{HumanEval 80\% 不是“80\% 编程能力”}
HumanEval 是约 164 个 Python function-completion tasks，以 unit tests 判断。分数受 prompt、sampling、pass@k、contamination 和 instruction tuning 影响，不能外推到 repository editing、debugging、安全或多文件软件工程。
\end{warningbox}

\subsection{Open code LLM landscape：provider、size 与 objective}

Meta Code Llama、BigCode StarCoder、DeepSeek Coder 等构成不同 family，另有 CodeGen、StableCode、CodeQwen、Granite。比较时应区分 base/instruct、context、FIM、license、languages、training date 与 benchmark protocol。

\lecturefigure{slide-58.jpg}{开放 code LLM landscape：Meta、BigCode、DeepSeek 与其他 families}{Loubna Ben Allal 官方 deck，第 58 页}

\begin{knowledgebox}{Code model 的首次术语消化}
\textbf{Base model} 主要学 next-token/FIM；\textbf{instruction-tuned model} 学指令回答；\textbf{MQA} 让所有 query heads 共享一组 KV；\textbf{GQA} 让若干 query heads 共享一组 KV，通常在质量与 KV-cache 之间折中。
\end{knowledgebox}

\teachervoice{00:55:00--00:57:20，Q\&A 中讲者说 code/text 训练总体相似，但 code assistant 更看重 long context、MQA/GQA inference efficiency 和 IDE-friendly smaller models。}

\subsection{BigCode：开放科学协作的目标}

BigCode 不只发布 weights，而是让 1,000+ 研究者/工程师参与 Slack，公开 data processing、training code、models 和 friendly license。这里的价值是能检查 decision trail，而不是把 collaboration size 当成 quality guarantee。

\lecturefigure{slide-59.jpg}{BigCode open-scientific collaboration：透明 data、processing、models 与社区}{Loubna Ben Allal 官方 deck，第 59 页}

\lecturefigure{slide-60.jpg}{BigCode ecosystem：The Stack、StarCoder 与社区 fine-tunes}{Loubna Ben Allal 官方 deck，第 60 页}

\begin{importantbox}{Ecosystem 的复用链}
The Stack 被多种 code models 使用；StarCoder checkpoints 又成为 StarChat、WizardCoder 等 fine-tunes 的 base。复用扩大影响，也意味着 upstream data quality、license 与 contamination decisions 会向下游传播。
\end{importantbox}

\teachervoice{00:40:00--00:43:10，讲者说 BigCode 发起的直接原因是 full data transparency：社区可以 inspect data、复用 processing code、检查模型，而不是只下载 weights。}

\subsection{本章小结}

Code LLM 生态从闭源 Copilot 快速扩展到公开 datasets、base models 和 fine-tunes。下一步问题不再是“有没有模型”，而是 release 是否真正开放、透明、可复现和负责。

